\chapter{Detalii de implementare}
\label{chapter:implementare}

În această secțiune se va discuta despre implementarea platformei, acoperind back-endul, front-endul, configurarea,
dezvoltarea și conectivitatea acestora. 

\section{Back-end}
\label{sec:implementare-backend}


Pentru back-end s-au analizat diferite tehnologii. Alegerea lor s-a bazat pe comparația descrisă mai jos.
Alegerea a fost influențată de mai mulți factori, printre care se numără performanța, ușurința de utilizare
și suportul extins pentru dezvoltarea Web . Go oferă un model de concurență eficient și o sintaxă simplă, 
ceea ce îl face potrivit pentru construirea rapidă de aplicații scalabile. Alt motiv principal a fost compatibilitatea
cu infrastructura aleasă pentru proiect, bazată pe Kubernetes și Docker, care se integrează foarte ușor cu Go.

Pe lângă Go, pentru serviciul de compilare a codului C++ s-a optat tot pentru consistență la Go, dar header-ul
care este inserat după detectarea structurilor de date pentru a crea raportări de status este scris in C++ 
pentru a asigura integrare nativă cu codul primit de la utilizator.


\begin{lstlisting}[language=Go, caption={Middleware Prometheus pentru monitorizarea duratei cererilor HTTP}]
template<typename Container>
class ContainerTracer : public Traceable 
public:
    ContainerTracer(const std::string& name, Container& container)
        : name(name), container(container) 
        globalTracer->registerObject(this);

class Traceable
public:
    virtual ~Traceable()
        globalTracer->unregisterObject(this);
\end{lstlisting}

Se poate observa că header-ul se bazează pe structuri generice din C++ pentru a permite orice tip de container
să fie integrat în acest tip general de raportare a operațiilor și conținutului.

Mai jos se află comparația limbajelor de care am discutat mai sus.

\begin{table}[h!]
\centering
\caption{Compararea limbajelor de programare care ar fi fost potrivite pentru VisioScience3D}
\begin{tabular}{|p{2.5cm}|p{4cm}|p{4cm}|p{4cm}|}
\hline
\textbf{Caracteristică} & \textbf{Go} & \textbf{Node.js (JavaScript)} & \textbf{Python} \\
\hline
\textbf{Performanță} & Compilat, foarte rapid, cu suport nativ pentru concurență eficientă (goroutine) & Interpretat, bun pentru I/O non-blocant, dar mai lent decât Go & Interpretat, mai lent, CPU-intensive tasks pot fi mai lente \\
\hline
\textbf{Concurență} & Goroutines ușor de folosit, eficiente din punct de vedere al resurselor & Model event-loop, asincron cu callback-uri/promises/async-await & Suport limitat (threading, asyncio), mai complex pentru concurență \\
\hline
\textbf{Simplitate și claritate} & Sintaxă simplă, tipuri de date clare, cod ușor de întreținut & Flexibil, dar poate duce la cod greu de urmărit (callback hell) & Sintaxă clară, dar structuri mari pot deveni greu de gestionat \\
\hline
\textbf{Ecosistem} & Ecosistem în creștere rapidă, multe librării pentru zona de microservicii, suport bun pentru REST, MongoDB & Ecosistem foarte matur și extins, multe librării, framework-uri & Ecosistem matur, multe librării pentru ML, backend, dar nu neapărat pentru performanță \\
\hline
\textbf{Scalabilitate} & Foarte scalabil, folosit în microservicii și sisteme distribuite & Scalabil pe I/O, dar single-threaded cu limite pe CPU & Scalabilitate limitată, se bazează pe external scaling și caching \\
\hline
\textbf{Deploy și containerizare} & Binar unic compilat, ușor de deployat în containere & Necesită runtime Node.js, imagini mai mari & Necesită runtime Python și librării, imagini relativ mari \\
\hline
\end{tabular}
\label{tab:comparatie-limbaje}
\end{table}


Fiecare microserviciu are propria configurație Kubernetes definită prin fișiere YAML pentru Deployment,
Service, ConfigMap și alte resurse necesare. Configurarea constă în containerizarea și automatizarea
livrării folosind Docker pentru containerizarea serviciilor și GitHub Actions pentru lanțul CI/CD,
care automatizează build-ul, testarea, și deploy-ul pe clusterul Kubernetes local.

Configurarea conexiunilor la baza de date MongoDB este centralizată în helperi, folosind driver-ul
oficial MongoDB pentru Go. Monitorizarea serviciilor este asigurată prin integrarea cu Prometheus,
cu metrici personalizate definite în cod și expuse la endpoint-uri REST.

\begin{lstlisting}[language=Go, caption={Middleware Prometheus pentru monitorizarea duratei cererilor HTTP}]
func prometheusMiddleware(next http.Handler) http.Handler {
    return http.HandlerFunc(func(w http.ResponseWriter, r *http.Request) {
        if r.URL.Path == "/feed/metrics" {
            next.ServeHTTP(w, r)
            return
        }
        start := prometheus.NewTimer(metrics.HTTPRequestDuration.WithLabelValues(r.Method, utils.NormalizePath(r.URL.Path)))
        next.ServeHTTP(w, r)
        start.ObserveDuration()
        metrics.HTTPRequestsTotal.WithLabelValues(r.Method, utils.NormalizePath(r.URL.Path), "200").Inc()
    })
}
\end{lstlisting}

\subsection{Dezvoltare back-end}
\label{subsec:implementare-backend}

Codul este organizat în pachete separate pentru rute, modele de date, pachet cu fișiere ajutătoare, metrici,
utils, panou grafana, pachet cu fișiere Kubernetes de configurare și middleware (ex: autentificare JWT).

Pentru rutare se utilizează biblioteca Gorilla Mux, iar fiecare endpoint implementează operații
CRUD pentru toate resursele din platformă precum quizuri, clase, molecule, formulare și altele.
Autentificarea și autorizarea utilizatorilor se realizează prin token-uri JWT, cu middleware dedicat.

Validarea datelor și gestionarea erorilor sunt tratate consistent, cu înregistrare de metrici
pentru succes și eșecuri, ceea ce permite monitorizarea calității serviciilor.

Pentru interacțiunea cu MongoDB, structurile Go sunt mapate la documente BSON, cu chei explicite 
și suport pentru ID-uri generate automat. Structurile BSON sunt convertite nativ in JSON pentru
serializare/deserializarea datelor și trimitea lor între servicii și front-end.

\begin{lstlisting}[language=Go, caption={Generare și validare token JWT}]
func GenerateToken(userID, role, email string) (string, error) {
    claims := CustomClaims{
        UserID: userID,
        Role:   role,
        Email:  email,
        RegisteredClaims: jwt.RegisteredClaims{
            ExpiresAt: jwt.NewNumericDate(time.Now().Add(24 * time.Hour)),
            IssuedAt:  jwt.NewNumericDate(time.Now()),
        },
    }
    token := jwt.NewWithClaims(jwt.SigningMethodHS256, claims)
    return token.SignedString(getJwtSecret())
}

func ParseToken(tokenString string) (*CustomClaims, error) {
    token, err := jwt.ParseWithClaims(tokenString, &CustomClaims{}, func(token *jwt.Token) (interface{}, error) {
        return getJwtSecret(), nil
    })
    if err != nil || !token.Valid {
        return nil, errors.New("invalid token")
    }
    return token.Claims.(*CustomClaims), nil
}
\end{lstlisting}

\subsection{Conectivitate}
\label{subsec:implementare-backend}

Microserviciile comunică prin API REST, iar traficul este expus și rulat în interiorul clusterului Kubernetes.
Pentru dezvoltare locală, se utilizează port-forwarding Kubernetes pentru a expune temporar serviciile
externe către mașina locală. Gestionarea rutelor externe și balansării numărului de cereri se realizează 
cu ajutorul Kong Ingress Controller, care asigură rutarea traficului HTTP către serviciile potrivite.

\begin{lstlisting}[language=Go, caption={Funcția CreateQuiz pentru crearea unui quiz nou în MongoDB}]
func CreateQuiz(w http.ResponseWriter, r *http.Request) {
    var input models.QuizInput
    if err := json.NewDecoder(r.Body).Decode(&input); err != nil {
        http.Error(w, err.Error(), http.StatusBadRequest)
        return
    }

    classID, err := primitive.ObjectIDFromHex(input.ClassID)
    if err != nil {
        http.Error(w, "Invalid class_id", http.StatusBadRequest)
        return
    }

    quiz := models.Quiz{
        ID:        primitive.NewObjectID(),
        Title:     input.Title,
        ClassID:   classID,
        Questions: input.Questions,
        CreatedAt: time.Now(),
    }

    collection := helpers.Client.Database("data-feed-db").Collection("quizzes")
    if _, err := collection.InsertOne(context.Background(), quiz); err != nil {
        http.Error(w, err.Error(), http.StatusInternalServerError)
        return
    }

    w.WriteHeader(http.StatusCreated)
    json.NewEncoder(w).Encode(quiz)
}
\end{lstlisting}

\begin{lstlisting}[language=Go, caption={Definirea rutelor și aplicarea middleware-ului JWT}]
r := mux.NewRouter()

r.Handle("/user/auth/register", handlers.Register).Methods("POST")
r.Handle("/user/auth/login", handlers.Login).Methods("POST")
r.Handle("/user/me", middleware.JWT(http.Handler(handlers.GetMe)))
r.Handle("/user/classes", middleware.JWT(http.Handler(handlers.CreateClass))).Methods("POST")
// Alte rute...

corsObj := gorillaHandlers.CORS(
    gorillaHandlers.AllowedOrigins([]string{"*"}),
    gorillaHandlers.AllowedMethods([]string{"GET", "POST", "PUT", "DELETE", "OPTIONS"}),
    gorillaHandlers.AllowedHeaders([]string{"Content-Type", "Authorization"}),
)

http.ListenAndServe(":8081", corsObj(r))
\end{lstlisting}


\section{Front-end}
\label{sec:implementare-frontend}

Pentru dezvoltarea front-end-ului am folosit React, împreună cu React Router DOM pentru navigare,
și tehnologia Three.js prin biblioteca React Three Fiber (R3F) și Drei pentru componente 3D reutilizabile.
Această combinație a permis crearea unei interfețe interactive plină de vizualizări 3D dinamice.

\subsection{Configurare front-end}
\label{subsec:implementare-frontend}

Mediul de dezvoltare a fost configurat cu Node.js și npm, folosind Vite ca bundler și server de dezvoltare rapidă.
În directorul proiectului, pachetele relevante instalate includ \texttt{react}, \texttt{react-router-dom}, 
\texttt{@react-three/fiber}, \texttt{@react-three/drei} și alte dependențe utile pentru animații și stilizare.
Tehnologiile menționate se îmbină cu folosirea Tailwind CSS pentru stilizare rapidă și eficientă,
permițând crearea de interfețe responsive. 

\subsection{Dezvoltare front-end}
\label{subsec:implementare-frontend}

Componenta principală a aplicației este structurată pe baza React și React Router pentru navigare între pagini.
Pentru vizualizările 3D s-au folosit componente R3F, cu suport Drei pentru facilități precum 
\texttt{useGLTF} pentru importul modelelor 3D, \texttt{useSpring} pentru animații și \texttt{Canvas}
pentru spațiul 3D.

Exemple de componente dezvoltate includ:
\begin{itemize}
  \item \textbf{Landing Pages} pentru materii (Astronomie, Chimie, Informatică) cu descrieri și elemente interactive.
  \item \textbf{Formulare} (ex: \texttt{ChemistryAddForm}) pentru adăugarea moleculelor cu 
  încărcare fișiere și validări.
  \item \textbf{Vizualizări 3D} (ex: \texttt{ChemistryBaloon}, \texttt{Drone}, \texttt{Island},
   \texttt{SpaceBackground}) animate și controlate interactiv.
  \item \textbf{Componente educaționale} pentru structuri de date și concepte de programare, 
  cu vizualizări grafice.
  \item \textbf{Dashboard-uri} și pagini pentru quiz-uri, clase, invitații, etc. cu apeluri către API-ul din backend.
\end{itemize}

Câteva exemple relevante de cod sunt prezentate mai jos:

Un exemplu de componentă standard React care utilizează Three.js pentru a crea un balon animat.
\begin{lstlisting}[language=JavaScript, basicstyle=\ttfamily\small, breaklines=true, caption={Exemplu de componentă React pentru un balon animat}]
const Baloon = (props) => {
  const [state, setState] = useState("idle");
  useEffect(() => {
    const interval = setInterval(() => {
      const states = ["idle","swayLeft","swayRight","tiltForward","tiltBackward","floatUp","floatDown"];
      setState(states[Math.floor(Math.random()*states.length)]);
    }, 2000);
    return () => clearInterval(interval);
  }, []);
  const { position, rotation } = useSpring({
    position: state==="floatUp" ? [0,0.3,0] : state==="floatDown" ? [0,-0.3,0] : [0,0,0],
    rotation: state==="swayLeft" ? [0,-0.1,0] : state==="swayRight" ? [0,0.1,0] : [0,0,0],
    config: { duration: 2000 }
  });
  return <a.mesh position={position} rotation={rotation} scale={props.scale}>...</a.mesh>;
};
\end{lstlisting}

\vspace{0.5em}

Alt exemplu este componenta meniului principal, care include o insulă complexă animată ce poate fi rotită cu mouse-ul.
\begin{lstlisting}[language=JavaScript, basicstyle=\ttfamily\small, breaklines=true, caption={Exemplu de animație a insulei meniului principal}]
const Island = ({ isRotating, setIsRotating }) => {
  const ref = useRef();
  const lastX = useRef(0);
  const alfa = 0.0075;
  const onDown = e => { setIsRotating(true); lastX.current = e.clientX || e.touches[0].clientX; };
  const onMove = e => {
    if (!isRotating) return;
    const currentX = e.clientX || e.touches[0].clientX;
    const delta = (currentX - lastX.current) / window.innerWidth;
    lastX.current = currentX;
    ref.current.rotation.y += delta * alfa * Math.PI;
  };
  const onUp = e => setIsRotating(false);
  useEffect(() => {
    window.addEventListener("mousedown", onDown);
    window.addEventListener("mouseup", onUp);
    window.addEventListener("mousemove", onMove);
    return () => { ... };
  }, [isRotating]);
  return <a.group ref={ref}>...</a.group>;
};
\end{lstlisting}

\subsection{Conectivitate}
\label{subsec:conectivitate-frontend}

Aplicația comunică cu backend-ul prin fetch API, efectuând cereri REST către serverul local pe portul 8000 datorită
mediului de dezvoltare în care s-a lucrat. S-au implementat următoarele pattern-uri:
\begin{itemize}
  \item Obținerea datelor dinamice (ex: formule, molecule, quiz-uri) în componente React folosind \texttt{useEffect}.
  \item Gestionarea token-ului JWT pentru autentificare, salvat în \texttt{localStorage} și folosit 
  în header-ul \texttt{Authorization}.
  \item Tratarea erorilor și afișarea mesajelor relevante utilizatorului.
  \item Formulare pentru inserare, ștergere și actualizare date pe server, cu răspunsuri și validări 
  vizuale în frontend.
\end{itemize}


Un exemplu de componentă React care face fetch pentru informațți pentru planeta Mercur se oservă mai jos.
\begin{lstlisting}[language=JavaScript, basicstyle=\ttfamily\small, breaklines=true, caption={Exemplu de fetch formule chimice cu loading și eroare}]
const MercuryFormulas = () => {
  const [formulas, setFormulas] = useState([]);
  const [loading, setLoading] = useState(true);
  const [error, setError] = useState("");
  useEffect(() => {
    fetch("http://localhost:8000/feed/shape/mercury")
      .then(res => {
        if(!res.ok) throw new Error();
        return res.json();
      })
      .then(data => setFormulas(Array.isArray(data)?data:[]))
      .catch(() => setError("Eroare la incarcare"))
      .finally(() => setLoading(false));
  }, []);
  return loading ? <p>Loading...</p> : error ? <p>{error}</p> : <ul>{formulas.map(f=><li key={f._id}>{f.formula.name}: {f.formula.expr}</li>)}</ul>;
};
\end{lstlisting}


Iar o secțiune importantă este vizualizarea 3D a moleculelor, care primește datele de la backend și crează o reprezentare
3D a atomilor și legăturilor chimice folosind Three.js.
\begin{lstlisting}[language=JavaScript, basicstyle=\ttfamily\small, breaklines=true, caption={Vizualizare 3D Molecule (atomi și legături cilindrice)}]
const Molecule = ({moleculeId}) => {
  const [atoms,setAtoms] = useState([]);
  const [bonds,setBonds] = useState([]);
  useEffect(() => {
    fetch(`/feed/chem/molecules/${moleculeId}/3d`)
      .then(r => r.json())
      .then(data => { setAtoms(data.atoms); setBonds(data.bonds); });
  }, [moleculeId]);
  useEffect(() => {
    if(groupRef.current) {
      const box = new THREE.Box3().setFromObject(groupRef.current);
      const center = new THREE.Vector3();
      box.getCenter(center);
      groupRef.current.position.sub(center);
    }
  }, [atoms,bonds]);
  return <group ref={groupRef}>
    {atoms.map((a,i)=><mesh key={i} position={[a.x,a.y,a.z]}><sphereGeometry args={[0.5,32,32]} /><meshStandardMaterial color="#ccc" /></mesh>)}
    {bonds.map((b,i)=> {
      const start=atoms[b.atom1], end=atoms[b.atom2];
      if(!start||!end) return null;
      const startV=new THREE.Vector3(start.x,start.y,start.z);
      const endV=new THREE.Vector3(end.x,end.y,end.z);
      const mid = startV.clone().add(endV).multiplyScalar(0.5);
      const dir = endV.clone().sub(startV).normalize();
      const len = startV.distanceTo(endV);
      const quat = new THREE.Quaternion().setFromUnitVectors(new THREE.Vector3(0,1,0), dir);
      return <mesh key={i} position={mid} quaternion={quat}><cylinderGeometry args={[0.1,0.1,len,16]} /><meshStandardMaterial color={0x888888} /></mesh>;
    })}
  </group>;
};
\end{lstlisting}


Iar crearea unei scene complexe ca cea din secțiunea de informatică pentru structura de date coadă cu priorități este realizată mai jos.

\begin{lstlisting}[language=JavaScript, basicstyle=\ttfamily\small, breaklines=true, caption={Exemplu de vizualizare de coadă cu priorități 3D folosind React Three Fiber}]
class PriorityQueue {
  constructor(comparator) {
    this.heap = [];
    this.compare = comparator;
  }
  push(val) { this.heap.push(val); this._siftUp(this.heap.length - 1); }
  pop() { if (this.heap.length === 0) return null; const top = this.heap[0]; const last = this.heap.pop(); if (this.heap.length > 0) { this.heap[0] = last; this._siftDown(0); } return top; }
  _shiftUp(idx) { /*  mutarea element in sus */ }
  _shiftDown(idx) { /* mutarea element in jos */ }
}

class TreeNode { constructor(value) { /* crearea unui nod din arbore */ } }
const buildTree = (arr) => { /* construirea arborelui dintr-un array */ };
const inorder = (node, arr = []) => { /* parcurgerea in ordine */ };
const computePositions = (node, x0, x1, y, gapY, list) => { /* calcularea pozitiilor nodurilor */ };

export default function PriorityQueueDemo() {
  const [type, setType] = useState("min");
  const [value, setValue] = useState("");
  const [pq, setPq] = useState(new PriorityQueue((a, b) => a - b));

  const handlePush = () => { if (value === "") return; pq.push(Number(value)); setValue(""); };
  const handlePop = () => { const top = pq.pop(); setPq(pq); };

  return (
    <div>
      <button onClick={handlePush}>push()</button>
      <button onClick={handlePop}>pop()</button>
      <div>{pq.heap.map((v, i) => <span key={i}>{v}</span>)}</div>
    </div>
  );
}
\end{lstlisting}

